% ===================================================================
%  図表第 8 号 — 収穫。— 狩。— 葡萄摘み。— 釣。
%
%  Ceci n'est PAS une transcription : c'est une traduction, faite en
%  2026, en japonais d'aujourd'hui. Voir texto/ja/10-zuhyo-01.tex
%  pour la consigne, qui vaut pour les seize fichiers.
% ===================================================================

%%K t08-noto-1 noto
\VUnotes{143.2mm}{%
\textit{(*) この語は十分に確かでないからである。亜麻を取るのか。葡萄を摘む
のか。林檎を取るのか。おそらく} \VUgras{mesono} \textit{を用いる方がよい。}
Mondo \textit{第十四巻二四二頁に提案されている。もっともこの新しい語根は今日
まで学会に採られておらず、イド語運動の常の規則に従って議されるのを待って
いる。}%
}

%%K t08-apar-1 apar
\VUsaut{5.0mm}
\VUtitre{15.0pt}{60}{図表第 8 号}
\VUsaut{3.0mm}
\VUcentre{13.2pt}{90}{\textit{第一場。}}
\VUsaut{3.0mm}
\VUpk{10.2pt}{40}{収穫 (*)。}
\VUsaut{4.0mm}
\VUpk{10.2pt}{40}{野の姿。}
\VUsaut{3.0mm}

%%K t08-c1-01-1 p
1. --- \VUgras{夏}が来て、野はまた姿を変えた。畝に託された種は芽を出し、
地から小さな緑の苗が伸びて穂をつけた。粒がふくらみ、やがて熟し、いまはもう
刈るばかりである。

%%K t08-c1-02-1 p
2. --- \VUgras{野の花}が咲いた。\VUgras{矢車菊} \textsuperscript{(1)}、
\VUgras{罌粟} \textsuperscript{(2)}、\VUgras{狐の手袋} \textsuperscript{(3)} が、
自分たちの新しい\VUgras{花}で麦畑の一色の地に明るい対を添えている。

%%K t08-c1-03-1 p
3. --- 果樹は葉を茂らせ、実は熟した。小さな\VUgras{桜桃の木} \textsuperscript{(4)}
は美しい\VUgras{桜桃} \textsuperscript{(5)} をたわわにつけ、子供たちがそれを摘んで
おいしそうに食べている。

%%K t08-c1-04-1 p
4. --- 羊の群はもう一日じゅう外にいられる。

%%K t08-c1-04-2 p
\VUgras{仔羊} \textsuperscript{(6)} は大きくなって、毛の厚い立派な\VUgras{牝羊}
\textsuperscript{(7)} と立派な\VUgras{牡羊} \textsuperscript{(8)} になった。静かに
\VUgras{牧場} \textsuperscript{(9)} の\VUgras{草}を食んでおり、その草には
\VUgras{雛菊}が点々と咲いている。

%%K t08-c1-04-3 p
やがてこの獣たちの\VUgras{羊毛}を刈ることになる。我々の衣の羅紗はそれで作る
のである。

%%K t08-tit-2 sub
\VUsaut{5.0mm}
\VUpk{10.2pt}{40}{牛飼。}
\VUsaut{3.4mm}

%%K t08-c1-05-1 p
5. --- \VUgras{牛飼} \textsuperscript{(10)} はあまり気にかけていないらしい。心に
あるのは地平を過ぎてゆく雷雨のことだけである。そして\VUgras{羊飼}
\textsuperscript{(13)} は、自分の\VUgras{水筒} \textsuperscript{(14)} を口へ運びながら、
もっと呑気そうに見える。

%%K t08-c1-06-1 p
6. --- 暑さは息苦しい。\VUgras{犬} \textsuperscript{(15)} までが涼みに来る。
\VUgras{小川} \textsuperscript{(16)} の澄んだ水を舐め、岸の草にうずくまっていた
気の毒な小さな\VUgras{蛙} \textsuperscript{(17)} を驚かせる。蛙は\VUgras{睡蓮}
\textsuperscript{(18)} の咲く澄んだ水へ跳びこむ。

%%K t08-c1-07-1 p
7. --- 一羽の\VUgras{鶺鴒} \textsuperscript{(19)} が小さな\VUgras{泉}
\textsuperscript{(20)} のそばで水を浴びている。その泉は二つの岩のあいだから湧き、
\VUgras{茨の茂み} \textsuperscript{(21)} と\VUgras{山査子の茂み} \textsuperscript{(22)}
の下へ隠れてゆく。

%%K t08-tit-3 sub
\VUsaut{5.0mm}
\VUpk{10.2pt}{40}{刈り入れ。}
\VUsaut{3.4mm}

%%K t08-c1-08-1 p
8. --- 村の子供たちには麦刈の時が面白くてならない。うれしそうに\VUgras{とんぼ
返り} \textsuperscript{(24)} をする。\VUgras{藁の山} \textsuperscript{(23)} の上で
である。\VUgras{刈る人} \textsuperscript{(26)} を手伝いもする。刈る人たちは
\VUgras{麦畑} \textsuperscript{(27)} を刈っている。使うのは\VUgras{大鎌}
\textsuperscript{(25)} で、\VUgras{砥石} \textsuperscript{(30)} でよく研いである。
子供たちは麦を集めて\VUgras{麦束} \textsuperscript{(31)} や\VUgras{小束}
\textsuperscript{(32)} にくくる。

%%K t08-c1-09-1 p
9. --- ときには年上の子が\VUgras{鎌} \textsuperscript{(12)} で、あの
\VUgras{黄金の茎} \textsuperscript{(28)} を刈ることを許される。その茎は実の詰まった
\VUgras{穂} \textsuperscript{(29)} を支えている。小さな子は\VUgras{家畜}
\textsuperscript{(33)} を見守りながら --- 家畜は\VUgras{草地} \textsuperscript{(34)} で
大きな\VUgras{菩提樹} \textsuperscript{(35)} の陰に草を食んでいる ---
自分たちの\VUgras{網} \textsuperscript{(36)} で美しい\VUgras{蝶}を捕まえる。

%%K t08-c1-09-2 p
何人かは\VUgras{荷車} \textsuperscript{(37)} にのぼって、\VUgras{草叉}
\textsuperscript{(38)} で差し上げられる麦束を並べさえする。

%%K t08-c1-10-1 p
10. --- \VUgras{平野} \textsuperscript{(39)} じゅう、\VUgras{雲雀}
\textsuperscript{(41)} が飛び立つあたりは活気に満ちている。誰もが刈り入れに
かかっている。しかし貧しい者が昔ながらの道具 --- \VUgras{大鎌、鎌}、
\VUgras{殻竿} \textsuperscript{(40)} --- を使っているのに対し、富んだ地主は自分の
麦や\VUgras{ライ麦} \textsuperscript{(43)} を\VUgras{刈取機} \textsuperscript{(42)}
で刈らせる。それから麦束を穀倉へ運ぶまでもなく、その場で大きな\VUgras{堆}
\textsuperscript{(44)} を積む。

%%K t08-c1-11-1 p
11. --- 次に\VUgras{脱穀機} \textsuperscript{(47)} を運んでくる。それを
\VUgras{移動汽罐} \textsuperscript{(45)} が\VUgras{伝動帯} \textsuperscript{(46)} で
回す。こうして粒は\VUgras{藁}から分かれる。蒔かれた畑を出ないままにである。

%%K t08-tit-4 sub
\VUsaut{5.0mm}
\VUpk{10.2pt}{40}{雷雨。}
\VUsaut{3.4mm}

%%K t08-c1-12-1 p
12. --- 刈り入れの時にはしばしば激しい\VUgras{雷雨} \textsuperscript{(53)} が
来る。まず遠くに\VUgras{雷}の音が聞える。強い\VUgras{西風} \textsuperscript{(54)}
が起こり、\VUgras{林} \textsuperscript{(49)} の一番大きな木を呻かせながらしならせる。
黒い\VUgras{雲} \textsuperscript{(56)} が空を攻め取る。それを\VUgras{稲妻}
\textsuperscript{(55)} の目もくらむ折れ線が引き裂く。\VUgras{雨}が、ときには
\VUgras{雹}までが降り出し、刈るばかりになった作物を倒してしまう。

%%K t08-c1-13-1 p
13. --- 雷はしばしば建物に火をつける。たとえば去年は\VUgras{風車小屋}
\textsuperscript{(50)} の屋根が灰になり、その\VUgras{翼} \textsuperscript{(51)} が
二枚砕けてしまった。

%%K t08-c1-14-1 p
14. --- 数時間ののち静けさが戻る。地平に輝かしい\VUgras{虹} \textsuperscript{(52)}
が現れ、しばし断たれた命を自然がまた取り上げる。\VUgras{小屋}
\textsuperscript{(48)} に逃げこんでいた村人は仕事に戻り、いま受けた損を直そうと
勇ましくかかる。

%%K t08-tit-5 sub
\VUsaut{6.0mm}
\VUcentre{13.2pt}{90}{\textit{第二場。}}
\VUsaut{3.6mm}

%%K t08-tit-6 sub
\VUpk{10.2pt}{40}{狩。--- 葡萄摘み。}
\VUsaut{1.4mm}
\VUpk{10.2pt}{40}{釣。}
\VUsaut{4.0mm}

%%K t08-c2-01-1 p
1. --- \VUgras{八月}の末か\VUgras{九月}の半ばに、土地によって違うが、狩の
季節が開く。日の初めの光が\VUgras{蜥蜴} \textsuperscript{(2)} を穴から誘い出す
やいなや、\VUgras{狩人} \textsuperscript{(5)} は自分の\VUgras{犬} \textsuperscript{(4)}
を連れて出かける。

%%K t08-c2-02-1 p
2. --- 厚い羅紗の丈夫な\VUgras{狩の衣} \textsuperscript{(9)} を着、革の
\VUgras{脚絆}をつける。\VUgras{蟇} \textsuperscript{(1)} の隠れる濡れた草地を歩く
ためである。\VUgras{猟銃} \textsuperscript{(6)} と\VUgras{獲物袋} \textsuperscript{(10)}
を取って、もう出てしまった。

%%K t08-c2-03-1 p
3. --- 狩の勢いが彼を葡萄畑の真中へ導く。そこでは葡萄摘みの最中である。
葡萄作りの子供たちは、いつもは道端で山羊の番をしているのに、今日は好きな所で
食ませておく。老いた\VUgras{牡山羊} \textsuperscript{(3)} を杭につないで ---
その獣は自由になればきっと逃げるからである --- 自分たちも楽しんでいる。一人は
\VUgras{葡萄籠} \textsuperscript{(12)} から房を一つ取って、\VUgras{汁}
\textsuperscript{(14)} を\VUgras{杯} \textsuperscript{(13)} に絞り、葡萄の汁（発酵
していない酒）を作って面白がっている。もう一人はいま編んだばかりの
\VUgras{葉の冠} \textsuperscript{(15)} を頭に載せている。そのそばで図々しい
\VUgras{雀} \textsuperscript{(16)} が搾り台のところまで葡萄を盗みに来ている。

%%K t08-c2-04-1 p
4. --- 子供が遊ぶあいだ、大人は働く。\VUgras{葡萄摘み} \textsuperscript{(18)}
たちは\VUgras{背負籠} \textsuperscript{(17)} で葡萄を穴倉へ運ぶ。一人の
\VUgras{桶屋} \textsuperscript{(19)} が\VUgras{酒樽} \textsuperscript{(20)} を直し、
\VUgras{側板} \textsuperscript{(22)} と\VUgras{鏡板} \textsuperscript{(21)} が
無事かを見、折れた\VUgras{たが} \textsuperscript{(23)} を替えている。あの大きな
\VUgras{槽} \textsuperscript{(25)} を作ったのもこの男で、そこに\VUgras{葡萄}
\textsuperscript{(26)} を入れて醸すのである。

%%K t08-c2-05-1 p
5. --- 醸しが終わると酒を抜き、残りを\VUgras{搾り台} \textsuperscript{(28)} に
載せる。大きな\VUgras{螺子} \textsuperscript{(29)} を少しずつ締めて汁を絞り、それが
\VUgras{受槽} \textsuperscript{(32)} に流れ、さらに\VUgras{酒} \textsuperscript{(31)}
を得るのである。

%%K t08-tit-8 sub
\VUsaut{6.0mm}
\VUpk{10.2pt}{40}{ビールと林檎酒。}
\VUsaut{3.4mm}

%%K t08-c2-06-1 p
6. --- \VUgras{酒倉} \textsuperscript{(27)} の壁に\VUgras{忽布} \textsuperscript{(30)}
の枝が這っている。ここではただの飾りの草にすぎないが、国によっては葡萄が
きわめて稀で、\VUgras{ビール}を醸すために忽布を育てる。

%%K t08-c2-07-1 p
7. --- 小さな\VUgras{林檎の木} \textsuperscript{(33)} は林檎をたわわにつけていて、
熟せばよい\VUgras{林檎酒}になる。自分の\VUgras{籠} \textsuperscript{(34)} の中で
\VUgras{金糸雀} \textsuperscript{(35)} と\VUgras{鶸} \textsuperscript{(36)} が、
外で自由に跳ねている雀を羨ましそうに眺めている。

%%K t08-c2-08-1 p
8. --- 見わたすかぎり丘の斜面には\VUgras{葡萄の杭} \textsuperscript{(40)} が
並んでいて、あの見事な\VUgras{葡萄の株} \textsuperscript{(39)} を支えている。株には
\VUgras{房}が垂れている。

%%K t08-c2-08-2 p
一年じゅう\VUgras{葡萄作り} \textsuperscript{(42)} は自分の\VUgras{葡萄畑}
\textsuperscript{(38)} の世話に働いてきた。いま美しい実りがその労に報いている。
\VUgras{葡萄摘みの女} \textsuperscript{(41)} たちが\VUgras{騾馬} \textsuperscript{(43)}
の引く車に載せた槽を満たし、誰もがうれしそうである。

%%K t08-tit-9 sub
\VUsaut{5.0mm}
\VUpk{10.2pt}{40}{釣。}
\VUsaut{3.4mm}

%%K t08-c2-09-1 p
9. --- \VUgras{釣人} \textsuperscript{(44)} たちも同じで、夜が明けるとすぐ自分の
\VUgras{釣竿} \textsuperscript{(45)} を持って水際に坐りに来ている。

%%K t08-c2-09-2 p
\VUgras{魚} \textsuperscript{(47)} は彼らが\VUgras{釣糸} \textsuperscript{(46)} を
水に入れるやいなや針にかかる。\VUgras{餌}を替える間もなく、糸の先で新しい
獲物がもがいている。

%%K t08-c2-09-3 p
それでも\VUgras{網打} \textsuperscript{(48)} の方がはるかに多く獲る。自分の
\VUgras{舟} \textsuperscript{(49)} から\VUgras{川} \textsuperscript{(50)} の真中へ
網を打つのである。

%%K t08-c2-10-1 p
10. --- 秋は出歩くによい季節である。歩いても、\VUgras{馬} \textsuperscript{(52)}
に乗っても、\VUgras{自転車} \textsuperscript{(53)} でも、\VUgras{自動車}
\textsuperscript{(54)} でもよい。\VUgras{道} \textsuperscript{(51)} は乾いていて、
人々は最後のよい日を惜しんで出かける。\VUgras{つばめ} \textsuperscript{(56)} は
もう屋根に集まり、暖かい国へ飛び立とうとしているからである。老いた
\VUgras{胡桃の木} \textsuperscript{(57)} の葉は黄ばみ、\VUgras{胡桃}が落ち、
\VUgras{木立} \textsuperscript{(58)} をなす高い\VUgras{木}も --- それは
\VUgras{高み} \textsuperscript{(59)} に立っている --- やがて裸になるであろう。

%%K t08-c2-11-1 p
11. --- \VUgras{太陽} \textsuperscript{(60)} は遅く昇り、昼は短くなり、
\VUgras{朝}にはもうかなりの寒さが感じられ、\VUgras{山鴫} \textsuperscript{(61)} の
群もすでに我々のもてなしのわるい土地を去ってゆく。
\VUsaut{6.0mm}
\VUfilet{22mm}
